\chapter{电磁辐射与辐射反作用}
\label{chap:maxwell-radiation}

给定局域电荷和电流以后，Maxwell 方程原则上已经决定所有时空点的场，但实验常常只测远区单位频率、单位立体角携带的能量。若每换一个源都重新解偏微分方程，这个可观测量很难直接读出；推迟 Green 函数把源到场的传播固定下来，Fourier 变换则把不同频率分开。

远处探测器看到的辐射并不是源在同一时刻的状态，而是守恒四流经过推迟 Green 函数传播后的横向部分。由推迟势取远区极限，再计算 Poynting 通量，可以得到适用于一般局域源的角分辨频谱。偶极辐射和加速点电荷都只是这个一般公式的特例；源受到辐射怎样反冲，则需要把场的能量动量重新反馈到源的运动方程中。

除特别说明外使用真空 SI 制。时间 Fourier 约定为
\begin{equation}
 \widetilde f(\omega)=\int_{-\infty}^{\infty}f(t)e^{\ii\omega t}\dd t,
 \qquad
 f(t)=\frac{1}{2\pi}\int_{-\infty}^{\infty}\widetilde f(\omega)e^{-\ii\omega t}\dd\omega .
 \label{eq:phys-rr-fourier-convention}
\end{equation}
于是 \(\widetilde{\dot f}=-\ii\omega\widetilde f\)。对实函数，\(\widetilde f(-\omega)=\widetilde f(\omega)^*\)，后文连续频谱均以 \(\omega>0\) 的正频率表示。对实矢量场，
\[
 \int_{-\infty}^{\infty}|\vb E(t)|^2\dd t
 =\frac1\pi\int_0^\infty|\widetilde{\vb E}(\omega)|^2\dd\omega .
\]
周期运动必须单独使用 Fourier 级数；连续变换中的 \(\delta\) 谱线不能当普通函数平方。观测方向的立体角元统一记为 \(\dd\Omega\)，周期源的基本角频率记为 \(\omega_0\)。

后文会出现多类彼此独立的小参数。固定 \(\omega>0\) 取远区极限时需要 \(a_{\rm src}/R\to0\)；若在有限探测距离使用平面波远区式，还需 \(kR\gg1\) 与 \(ka_{\rm src}^2/R\ll1\)。对空间上局域的有限尺寸源，若电流在所考虑的时间区间内具有紧支撑，可取 \(a_{\rm src}\) 使 \(J^\mu(\vb x,t)=0\) 于 \(|\vb x|>a_{\rm src}\)；对快速衰减但非紧支撑的源，则用相应空间矩的特征尺度代替，并必须另外控制尾部余项。统一定义结构长波参数
\begin{equation}
 \kappa_{\rm src}:=\frac{\omega a_{\rm src}}{c}.
 \label{eq:phys-rr-source-size-control}
\end{equation}
它既控制给定源的空间多极展开，也控制平稳有限尺寸自作用核的低频结构展开；同一个物理尺度不再另造第二个符号。散射的软频展开则由 \(\epsilon_{\rm soft}:=\omega T_{\rm sc}\) 控制，辐射反作用降阶的外场梯度部分由 \(\epsilon_{\rm grad}:=\tau_{\rm rr}/T\) 控制。协变 Landau--Lifshitz 近似还要另行检查强场加速度参数。不同小参数不能互相替代：例如 \(\omega T_{\rm sc}\ll1\) 并不推出 \(\kappa_{\rm src}\ll1\)。

\section{从 Maxwell 方程到一般频率的精确辐射谱}

本节始终先固定 \(\omega>0\)，再取 \(R\to\infty\) 的 \(R^{-1}\) 辐射系数；只有母式建立以后才讨论 \(\omega\to0\)。因此“软频”“长波”和“周期”都是不同方向的后续特殊化，而不是一般辐射公式的前提。

设局域四流 \(J^\mu=(c\rho_{\rm e},\vb J)\) 满足
\begin{equation}
 \partial_t\rho_{\rm e}+\nabla\cdot\vb J=0.
 \label{eq:phys-rr-continuity}
\end{equation}
对观测方向 \(\vb n\)（\(|\vb n|=1\)）定义光锥上的时空 Fourier 振幅
\begin{align}
 \mathcal C_{\rm e}(\omega,\vb n)
 &:=\int\dd t\dd^3x\,
 \rho_{\rm e}(\vb x,t)e^{\ii\omega(t-\vb n\cdot\vb x/c)},\\
 \boldsymbol{\mathcal J}(\omega,\vb n)
 &:=\int\dd t\dd^3x\,
 \vb J(\vb x,t)e^{\ii\omega(t-\vb n\cdot\vb x/c)}.
 \label{eq:phys-rr-exact-current-amplitude}
\end{align}
把辐射壳上的四流振幅记为
\[
 \mathcal J^\mu(\omega,\vb n):=
 \bigl(c\mathcal C_{\rm e}(\omega,\vb n),\boldsymbol{\mathcal J}(\omega,\vb n)\bigr).
\]
连续性方程在光子四动量 \(k^\mu=(\omega/c,(\omega/c)\vb n)\) 上给出
\begin{equation}
 \vb n\cdot\boldsymbol{\mathcal J}=c\mathcal C_{\rm e},
 \label{eq:phys-rr-on-shell-current-conservation}
\end{equation}
即 \(k_\mu\mathcal J^\mu=0\)。因此远区电场只依赖电流的横向投影。令
\[
 \Pi^{\perp}_{ij}(\vb n)=\delta_{ij}-n_in_j,
 \qquad
 \vb{\Pi}_\perp(\vb n):=(\Pi^{\perp}_{ij})_{i,j=1}^{3},
\]
则固定 \(\omega\) 的精确远区振幅为
\begin{equation}
 \widetilde{\vb E}_{\rm rad}(R\vb n,\omega)
 =\frac{\ii\mu_0\omega e^{\ii\omega R/c}}{4\pi R}
 \vb{\Pi}_{\perp}(\vb n)\cdot\boldsymbol{\mathcal J}(\omega,\vb n)
 +O(R^{-2}).
 \label{eq:phys-rr-exact-far-field-current}
\end{equation}
\eqnrefs{eq:phys-rr-exact-far-field-current} 对任意亚光速运动、任意源尺寸参数 \(\omega a_{\rm src}/c\) 和非周期源都成立；低速、长波与周期性只在后续特殊化时施加。正频率角分辨能谱因此为
\begin{equation}
 \frac{\dd^2E}{\dd\omega\dd\Omega}
 =\frac{\mu_0\omega^2}{16\pi^3c}
 \Pi^{\perp}_{ij}\mathcal J_i\mathcal J_j^*
 =\frac{\omega^2}{16\pi^3\varepsilon_0c^3}
 \left|\vb{\Pi}_{\perp}\cdot\boldsymbol{\mathcal J}\right|^2.
 \label{eq:phys-rr-exact-radiation-spectrum-current}
\end{equation}
这是给定守恒源在一般频率下的辐射公式。非周期且总辐射能有限的过程可直接使用它；无限持续的周期源则应改用周期平均功率。

对角度积分，可把源的空间相干性写成一个不依赖具体几何的双点核。令 \(\Delta\vb x=\vb x-\vb x'\)、\(r=|\Delta\vb x|\)、\(\widehat{\vb r}=\Delta\vb x/r\)、\(z=\omega r/c\)，记 \(j_\ell(z)\) 为第一类球 Bessel 函数。定义
\begin{equation}
 \mathcal K_{ij}(z,\widehat{\vb r})
 :=\int\dd\Omega\,\Pi^{\perp}_{ij}(\vb n)e^{-\ii z\vb n\cdot\widehat{\vb r}}
 =\frac{4\pi}{3}
 \left[(2j_0-j_2)\delta_{ij}+3j_2\widehat r_i\widehat r_j\right].
 \label{eq:phys-rr-transverse-coherence-kernel}
\end{equation}
于是
\begin{equation}
 \frac{\dd E}{\dd\omega}
 =\frac{\mu_0\omega^2}{16\pi^3c}
 \int\dd t\dd t'\dd^3x\dd^3x'\,
 e^{\ii\omega(t-t')}
 J_i(\vb x,t)J_j(\vb x',t')
 \mathcal K_{ij}\!\left(\frac{\omega r}{c},\widehat{\vb r}\right).
 \label{eq:phys-rr-exact-angle-integrated-current-spectrum}
\end{equation}
当 \(\omega r/c\ll1\) 时核近似同相；当 \(\omega r/c\gtrsim1\) 时，不同位置之间出现振荡相消。空间长波展开因此只是这个精确相干核的低频空间极限。

\begin{figure}[htbp]
\centering
\begin{adjustbox}{max width=0.90\linewidth}
\begin{tikzpicture}[x=1.25cm,y=1.05cm]
  \draw[lecturearrow] (0,0) -- (6.4,0)
    node[right,lecturelabel] {$\kappa_{\rm src}=\omega a_{\rm src}/c$};
  \draw[lecturearrow] (0,0) -- (0,5.2)
    node[above,lecturelabel] {$\epsilon_{\rm soft}=\omega T_{\rm sc}$};

  \draw[dashed,line width=0.65pt] (2.85,0) -- (2.85,4.75);
  \draw[dashed,line width=0.65pt] (0,1.35) -- (5.95,1.35);

  \node[lecturelabel] at (1.35,3.75)
    {long-wave region\\$\kappa_{\rm src}\ll1$\\multipole expansion};
  \node[lecturelabel] at (4.45,0.64)
    {soft-frequency region\\$\epsilon_{\rm soft}\ll1$\\asymptotic-state control};
  \node[lecturelabel] at (1.35,0.63)
    {overlap\\soft multipole limit};
  \node[lecturelabel] at (4.45,3.25)
    {general-frequency region\\full spatial phase\\and trajectory history};

  \node[draw,rounded corners,fill=white,lecturelabel,inner sep=4pt] at (4.45,4.45)
    {exact parent formula:\\
     $\displaystyle \frac{\dd^2E}{\dd\omega\dd\Omega}
     =\frac{\omega^2}{16\pi^3\varepsilon_0c^3}
     \bigl|\vb{\Pi}_\perp\!\cdot\!\vb{\mathcal J}\bigr|^2$};

  \node[below,lecturelabel,text width=7.1cm] at (3.1,-0.45)
    {$a_{\rm src}/c$ measures spatial retardation, whereas
     $T_{\rm sc}$ measures source-evolution time; neither small
     parameter implies the other.};
\end{tikzpicture}

\end{adjustbox}
\caption{时间软展开与空间长波展开是两个独立方向。\(\epsilon_{\rm soft}=\omega T_{\rm sc}\) 控制散射历史的软频极限，\(\kappa_{\rm src}=\omega a_{\rm src}/c\) 控制源内空间相位的多极展开；整个参数平面都由精确母\eqnrefs{eq:phys-rr-exact-radiation-spectrum-current} 覆盖。}
\label{fig:phys-rr-frequency-regime-map}
\end{figure}

\begin{exbox}
有限尺寸相干性可以直接从母式计算。若球对称静态分布 \(\rho_{{\rm e},0}(\vb x)\) 作小振幅、非相对论刚性平移（\(|\dot{\vb X}|/c\ll1\)），精确密度是 \(\rho_{\rm e}(\vb x,t)=\rho_{{\rm e},0}(\vb x-\vb X(t))\)。在线性阶，\(\delta\rho_{\rm e}=-\vb X\cdot\nabla\rho_{{\rm e},0}\)、\(\vb J=\rho_{{\rm e},0}\dot{\vb X}\)，二者逐阶满足连续性方程。于是
\[
 \boldsymbol{\mathcal J}(\omega,\vb n)
 =\widetilde{\dot{\vb X}}(\omega)
 \widetilde\rho_{{\rm e},0}\!\left(\frac{\omega}{c}\vb n\right).
\]
若总电荷为 \(Q\)，取半径 \(a_{\rm src}\) 的均匀薄球壳
\[
 \rho_{{\rm e},{\rm shell}}(r)=\frac{Q}{4\pi a_{\rm src}^2}\delta(r-a_{\rm src})
\]
以及以 \(a_{\rm src}\) 为径向尺度的高斯云
\[
 \rho_{{\rm e},{\rm G}}(r)=\frac{Q}{\pi^{3/2}a_{\rm src}^3}e^{-r^2/a_{\rm src}^2}.
\]
二者的归一化 Fourier 形状因子分别为 \(\sin\kappa_{\rm src}/\kappa_{\rm src}\) 与 \(e^{-\kappa_{\rm src}^2/4}\)，故强度比为
\[
 \frac{I_{\rm shell}}{I_{\rm point}}
 =\left(\frac{\sin\kappa_{\rm src}}{\kappa_{\rm src}}\right)^2,
 \qquad
 \frac{I_{\rm G}}{I_{\rm point}}=e^{-\kappa_{\rm src}^2/2},
 \qquad \kappa_{\rm src}=\frac{\omega a_{\rm src}}{c}.
\]
二者在 \(\kappa_{\rm src}\ll1\) 时都趋于点源，但离开长波区后的振荡与衰减依赖完整电荷分布，而不存在普适的单一有限尺寸修正系数。
\end{exbox}

\begin{proof}[解]
Lorenz 规范下的势满足
\[
 \left(\nabla^2-c^{-2}\partial_t^2\right)\Phi=-\rho_{\rm e}/\varepsilon_0,
 \qquad
 \left(\nabla^2-c^{-2}\partial_t^2\right)\vb A=-\mu_0\vb J .
\]
对时间作 Fourier 变换，并用出射 Helmholtz 核
\(G_k(\vb R-\vb x)=e^{\ii k|\vb R-\vb x|}/(4\pi|\vb R-\vb x|)\)。
在固定 \(\omega>0\) 的远区，
\[
 G_k(R\vb n-\vb x)
 =\frac{e^{\ii kR}}{4\pi R}e^{-\ii k\vb n\cdot\vb x}
 +O(R^{-2}),
\]
故
\[
 \widetilde\Phi_{\rm rad}
 =\frac{e^{\ii kR}}{4\pi\varepsilon_0R}\mathcal C_{\rm e},
 \qquad
 \widetilde{\vb A}_{\rm rad}
 =\frac{\mu_0e^{\ii kR}}{4\pi R}\boldsymbol{\mathcal J}.
\]
连续性方程在同一质壳 Fourier 分量上给出
\(-\ii\omega\mathcal C_{\rm e}+\ii k\vb n\cdot\boldsymbol{\mathcal J}=0\)。
因此 \(\widetilde{\vb E}=\ii\omega\widetilde{\vb A}-\nabla\widetilde\Phi\)
中的纵向部分抵消，只剩
\[
 \widetilde{\vb E}_{\rm rad}
 =\frac{\ii\mu_0\omega e^{\ii kR}}{4\pi R}
 \vb\Pi_\perp\cdot\boldsymbol{\mathcal J}+O(R^{-2}).
\]
再用 \(\widetilde{\vb B}_{\rm rad}=\vb n\times\widetilde{\vb E}_{\rm rad}/c\)
和正频率 Parseval 关系
\( \dd^2E/(\dd\omega\dd\Omega)=R^2|\widetilde{\vb E}_{\rm rad}|^2/(\pi\mu_0c)\)，
即得\eqnrefs{eq:phys-rr-exact-radiation-spectrum-current}。
\end{proof}

\begin{proof}[解]
取 \(\widehat{\vb r}\) 为极轴。标量角积分为
\[
 \int\dd\Omega\,e^{-\ii z\vb n\cdot\widehat{\vb r}}=4\pi j_0(z).
\]
旋转对称性把二阶角矩锁成 \(I_{ij}=A(z)\delta_{ij}+B(z)\hat r_i\hat r_j\)。
\[
 I_{ii}=\int\dd\Omega\,(\vb n\cdot\vb n)e^{-\ii z\vb n\cdot\hat{\vb r}}=\int\dd\Omega\,e^{-\ii z\vb n\cdot\hat{\vb r}}=4\pi j_0(z),
\]
故 \(3A+B=4\pi j_0\)。再沿极轴取分量。记 \(\mu=\vb n\cdot\hat{\vb r}\)，注意到
\[
 \frac{\dd^2}{\dd z^2}e^{-\ii z\mu}=(-\ii\mu)^2e^{-\ii z\mu}=-\mu^2 e^{-\ii z\mu},
\]
于是
\[
 I_{rr}=\hat r_i\hat r_j I_{ij}=\int\dd\Omega\,\mu^2 e^{-\ii z\mu}
 =-\frac{\dd^2}{\dd z^2}\int\dd\Omega\,e^{-\ii z\mu}=-4\pi j_0''(z).
\]
球 Bessel 递推 \(j_0'=-j_1\)、\(j_1'=j_0-\frac{2}{z}j_1\) 与 \(j_0+j_2=\frac{3}{z}j_1\) 联立得 \(j_0''=\frac13(2j_2-j_0)\)，故 \(I_{rr}=A+B=\frac{4\pi}{3}(j_0-2j_2)\)。与 \(3A+B=4\pi j_0\) 联立解出 \(A=\frac{4\pi}{3}(j_0+j_2)\)、\(B=-4\pi j_2\)，即
\[
 \int\dd\Omega\,n_in_j e^{-\ii z\vb n\cdot\hat{\vb r}}
 =\frac{4\pi}{3}\left[(j_0+j_2)\delta_{ij}
 -3j_2\hat r_i\hat r_j\right].
\]
从标量积分中减去这一式，便得到
\[
 \int\dd\Omega\,\Pi^\perp_{ij}e^{-\ii z\vb n\cdot\widehat{\vb r}}
 =\frac{4\pi}{3}\left[(2j_0-j_2)\delta_{ij}
 +3j_2\widehat r_i\widehat r_j\right],
\]
即\eqnrefs{eq:phys-rr-transverse-coherence-kernel}。把
\(\mathcal J_i\mathcal J_j^*\) 写成两份源积分后先做角积分，直接得到
\eqnrefs{eq:phys-rr-exact-angle-integrated-current-spectrum}。
\end{proof}

点电荷世界线
\[
 \vb J(\vb x,t)=\sum_aq_ac\boldsymbol\beta_a(t)
 \delta^{(3)}(\vb x-\vb r_a(t)),
 \qquad
 \varphi_a=t-\frac{\vb n\cdot\vb r_a}{c}
\]
其中 $\boldsymbol\beta_a=\vb v_a/c$ 是粒子 $a$ 的速度矢量（以 $c$ 为单位，三维速度空间中的矢量），$\vb r_a(t)$ 为其世界线，$\varphi_a$ 为推迟时刻。
代入母式后给出两种等价表示：
\begin{align}
 \frac{\dd^2E}{\dd\omega\dd\Omega}
 &=\frac{\omega^2}{16\pi^3\varepsilon_0c}
 \left|\sum_aq_a\int\dd t\,
 \vb\Pi_\perp\cdot\boldsymbol\beta_a e^{\ii\omega\varphi_a}\right|^2,
 \label{eq:phys-rr-point-spectrum-velocity}\\
 \frac{\dd^2E}{\dd\omega\dd\Omega}
 &=\frac{1}{16\pi^3\varepsilon_0c}
 \left|\sum_aq_a\int\dd t\,
 \frac{\vb n\times[(\vb n-\boldsymbol\beta_a)\times
 \dot{\boldsymbol\beta}_a]}
 {(1-\vb n\cdot\boldsymbol\beta_a)^2}
 e^{\ii\omega\varphi_a}\right|^2.
 \label{eq:phys-rr-exact-point-spectrum}
\end{align}
第一式适合直接代入轨道；第二式把辐射来源明确限制为加速度，并在渐近匀速散射中具有更好的收敛性。二者对有限正频率都不要求非相对论或长波近似。

\begin{proof}[解]
对点流，空间积分先给出
\[
 \boldsymbol{\mathcal J}
 =c\sum_aq_a\int\dd t\,\boldsymbol\beta_a e^{\ii\omega\varphi_a},
 \qquad
 \dot\varphi_a=1-\vb n\cdot\boldsymbol\beta_a .
\]
定义
\[
 \vb g_a=\frac{\vb n\times(\vb n\times\boldsymbol\beta_a)}
 {1-\vb n\cdot\boldsymbol\beta_a}
 =-\frac{\vb\Pi_\perp\cdot\boldsymbol\beta_a}
 {1-\vb n\cdot\boldsymbol\beta_a}.
\]
在渐近匀速轨道上用绝热因子消去边界项，分部积分得到
\[
 \ii\omega\int\dd t\,
 \vb\Pi_\perp\cdot\boldsymbol\beta_a e^{\ii\omega\varphi_a}
 =\int\dd t\,\dot{\vb g}_a e^{\ii\omega\varphi_a}.
\]
对 \(\vb g_a\) 求导，商法则与三重叉积恒等式给出
\[
 \dot{\vb g}_a=
 \frac{\vb n\times[(\vb n-\boldsymbol\beta_a)\times
 \dot{\boldsymbol\beta}_a]}
 {(1-\vb n\cdot\boldsymbol\beta_a)^2}.
\]
代回速度形式后，振幅中的 \(1/\omega\) 与谱前因子中的 \(\omega^2\)
相消，得到\eqnrefs{eq:phys-rr-exact-point-spectrum}。这也直接说明匀速段
\(\dot{\boldsymbol\beta}_a=0\) 不贡献有限频率辐射。
\end{proof}

周期性只离散时间频率。若
\(\vb J(\vb x,t+T)=\vb J(\vb x,t)\)、\(\omega_N=N\omega_0\)，定义
\begin{equation}
 \boldsymbol{\mathcal J}_N(\vb n)
 =\int\dd^3x\,\vb J_N(\vb x)
 e^{-\ii(\omega_N/c)\vb n\cdot\vb x},
 \qquad
 \vb J_N=\frac1T\int_0^T\vb J e^{\ii N\omega_0t}\dd t .
 \label{eq:phys-rr-periodic-current-amplitude}
\end{equation}
第 \(N\) 条谱线的精确平均功率为
\begin{equation}
 \frac{\dd P_N}{\dd\Omega}
 =\frac{\mu_0\omega_N^2}{8\pi^2c}
 \left|\vb\Pi_\perp\cdot\boldsymbol{\mathcal J}_N\right|^2 .
 \label{eq:phys-rr-exact-periodic-line-power}
\end{equation}
只有在得到平均功率后才写离散 \(\delta\) 谱，从而避免平方 \(\delta\) 分布。

空间多极展开则另需
\begin{equation}
 \kappa_{\rm src}=\frac{\omega a_{\rm src}}{c}\ll1 .
 \label{eq:phys-rr-multipole-control}
\end{equation}
定义
\begin{align}
 \vb p&=\int\rho_{\rm e}\vb x\,\dd^3x,
 &\vb m&=\frac12\int\vb x\times\vb J\,\dd^3x,
 \label{eq:phys-rr-dipole-defs}\\
 \mathcal Q^{\rm E}_{ij}
 &=\int\rho_{\rm e}(3x_ix_j-r^2\delta_{ij})\dd^3x,
 &\mathcal Q^{\rm E}_{ii}&=0 ,
 \label{eq:phys-rr-quadrupole-def}
\end{align}
并记 \(I_{ij}=\int\rho_{\rm e}x_ix_j\dd^3x\)。连续性方程给出
\begin{equation}
 \int x_jJ_i\dd^3x=\frac12\dot I_{ij}-\epsilon_{ijk}m_k .
 \label{eq:phys-rr-first-current-moment}
\end{equation}
于是
\begin{equation}
 \vb\Pi_\perp\cdot\boldsymbol{\mathcal J}
 =-\ii\omega\vb\Pi_\perp\cdot\widetilde{\vb p}
 +\ii k\,\vb n\times\widetilde{\vb m}
 -\frac{k\omega}{6}\vb\Pi_\perp\cdot
 (\widetilde{\vb{\mathcal Q}}^{\rm E}\cdot\vb n)
 +\boldsymbol{\mathcal R}_2 .
 \label{eq:phys-rr-multipole-current-amplitude}
\end{equation}
余项从更高空间矩开始，量级由 \(\kappa_{\rm src}\) 控制。逆 Fourier 变换为
\begin{equation}
 \vb E_{\rm rad}
 =\frac{1}{4\pi\varepsilon_0c^2R}
 \left\{\vb n\times(\vb n\times\ddot{\vb p})
 +\frac1c\vb n\times\ddot{\vb m}
 +\frac1{6c}\vb n\times[
 \vb n\times(\dddot{\vb{\mathcal Q}}^{\rm E}\cdot\vb n)]
 \right\}_{t-R/c}+\vb E_{\rm rem}.
 \label{eq:phys-rr-multipole-far-field}
\end{equation}
完整角积分后，
\begin{align}
 \frac{\dd E_{E1}}{\dd\omega}
 &=\frac{\omega^4}{6\pi^2\varepsilon_0c^3}
 \widetilde p_i\widetilde p_i^*,
 \label{eq:phys-rr-e1-spectrum}\\
 \frac{\dd E_{M1}}{\dd\omega}
 &=\frac{\omega^4}{6\pi^2\varepsilon_0c^5}
 \widetilde m_i\widetilde m_i^*,
 \label{eq:phys-rr-m1-spectrum}\\
 \frac{\dd E_{E2}}{\dd\omega}
 &=\frac{\omega^6}{720\pi^2\varepsilon_0c^5}
 \widetilde{\mathcal Q}^{\rm E}_{ij}
 (\widetilde{\mathcal Q}^{\rm E}_{ij})^* .
 \label{eq:phys-rr-e2-spectrum}
\end{align}
对应的瞬时功率为
\begin{equation}
 P_{E1}=\frac{|\ddot{\vb p}|^2}{6\pi\varepsilon_0c^3},
 \quad
 P_{M1}=\frac{|\ddot{\vb m}|^2}{6\pi\varepsilon_0c^5},
 \quad
 P_{E2}=\frac{\dddot{\mathcal Q}^{\rm E}_{ij}
 \dddot{\mathcal Q}^{\rm E}_{ij}}{720\pi\varepsilon_0c^5}.
 \label{eq:phys-rr-multipole-time-powers}
\end{equation}
若源同时周期化，则
\begin{align}
 P_{E1,N}&=\frac{(N\omega_0)^4}{3\pi\varepsilon_0c^3}p_{i,N}p_{i,N}^*,\\
 P_{M1,N}&=\frac{(N\omega_0)^4}{3\pi\varepsilon_0c^5}m_{i,N}m_{i,N}^*,\\
 P_{E2,N}&=\frac{(N\omega_0)^6}{360\pi\varepsilon_0c^5}
 \mathcal Q^{\rm E}_{ij,N}(\mathcal Q^{\rm E}_{ij,N})^* .
 \label{eq:phys-rr-periodic-multipole-lines}
\end{align}
周期化与长波展开彼此独立；两者可交换次序，但必须分别满足各自条件。

散射的软频极限由
\begin{equation}
 \epsilon_{\rm soft}=\omega T_{\rm sc}\ll1
 \label{eq:phys-rr-soft-control}
\end{equation}
控制。对入射、出射外线分别取 \(\varsigma_A=-1,+1\)，领先软因子为
\begin{equation}
 \boldsymbol{\mathcal S}^{(0)}(\vb n)
 =\sum_A\varsigma_Aq_A
 \frac{\vb n\times(\vb n\times\boldsymbol\beta_A)}
 {1-\vb n\cdot\boldsymbol\beta_A},
 \label{eq:phys-rr-leading-soft-factor}
\end{equation}
并给出
\begin{equation}
 \left.\frac{\dd^2E}{\dd\omega\dd\Omega}\right|_{\omega\to0}
 =\frac{|\boldsymbol{\mathcal S}^{(0)}|^2}
 {16\pi^3\varepsilon_0c}.
 \label{eq:phys-rr-soft-angular-spectrum}
\end{equation}
若 \(\sum_A\varsigma_Aq_A=0\)，令
\(\gamma_A=1/\sqrt{1-\beta_A^2}\)（\(\beta_A=|\boldsymbol\beta_A|\)）、\(\Gamma_{AB}=\gamma_A\gamma_B(1-\boldsymbol\beta_A\cdot\boldsymbol\beta_B)\)、
\(\beta_{AB}=\sqrt{1-\Gamma_{AB}^{-2}}\)，角积分为
\begin{equation}
 \left.\frac{\dd E}{\dd\omega}\right|_{\omega\to0}
 =-\frac1{4\pi^2\varepsilon_0c}
 \sum_{A,B}\varsigma_A\varsigma_Bq_Aq_B
 \left[\frac{\operatorname{artanh}\beta_{AB}}{\beta_{AB}}-1\right].
 \label{eq:phys-rr-soft-integrated-spectrum}
\end{equation}
单电荷非相对论偏转以及与偶极路线的重叠极限分别是
\begin{align}
 \left.\frac{\dd E}{\dd\omega}\right|_{\omega\to0}
 &=\frac{q^2|\Delta\vb v|^2}{6\pi^2\varepsilon_0c^3}+O(c^{-5}),
 \label{eq:phys-rr-soft-nonrel-limit}\\
 \left.\frac{\dd E_{E1}}{\dd\omega}\right|_{\omega\to0}
 &=\frac{|\Delta\dot{\vb p}|^2}{6\pi^2\varepsilon_0c^3}.
 \label{eq:phys-rr-e1-soft-limit}
\end{align}
零频振幅的时域形式是电磁记忆：
\begin{equation}
 \int_{-\infty}^{\infty}\vb E_{\rm rad}(R\vb n,u)\dd u
 =\frac{\boldsymbol{\mathcal S}^{(0)}(\vb n)}
 {4\pi\varepsilon_0cR}.
 \label{eq:phys-rr-em-memory-soft-relation}
\end{equation}
这些式子只使用时间软极限，不应与固定尺寸源的长波多极展开混为一谈。

\begin{proof}[解]
在点世界线加速度振幅中使用前述 \(\vb g_a\)。当
\(\omega T_{\rm sc}\ll1\) 时，
\[
 \int\dd t\,\dot{\vb g}_a e^{\ii\omega\varphi_a}
 =\vb g_a(+\infty)-\vb g_a(-\infty)+O(\epsilon_{\rm soft}),
\]
相加即得\eqnrefs{eq:phys-rr-leading-soft-factor}。角积分利用
\[
 \frac1{AB}=\int_0^1\frac{\dd x}{[xA+(1-x)B]^2},
 \qquad
 \int\frac{\dd\Omega}{(1-\vb n\cdot\vb w)^2}
 =\frac{4\pi}{1-w^2},
\]
其中 \(\vb w=\vb v/c\) 为辐射电荷的速度（以 \(c\) 为单位），\(w^2=|\vb w|^2\)。
再由电荷守恒减去与
\((\sum_A\varsigma_Aq_A)^2\) 成正比的常数项，得到
\eqnrefs{eq:phys-rr-soft-integrated-spectrum}。
\end{proof}

\begin{figure}[htbp]
\centering
\includegraphics[width=0.68\linewidth]{form_factor_shell.pdf}
\caption{精确有限尺寸形状因子与长波展开。薄球壳的精确强度比为 \(\sinc^2\kappa_{\rm src}\)，高斯云为 \(e^{-\kappa_{\rm src}^2/2}\)；虚线是各自的低 \(\kappa_{\rm src}\) 展开。同一个控制参数 \(\kappa_{\rm src}=\omega a_{\rm src}/c\) 并不意味着离开长波区后具有相同修正。}
\label{fig:phys-rr-form-factor-shell}
\end{figure}

\begin{proof}[解]
在 \(\kappa_{\rm src}\ll1\) 时展开空间相位，
\[
 e^{-\ii k\vb n\cdot\vb x}
 =1-\ii k n_jx_j+O(\kappa_{\rm src}^2).
\]
连续性方程分部积分给出
\[
 \int J_i\dd^3x=\dot p_i,\qquad
 \int x_jJ_i\dd^3x=\frac12\dot I_{ij}-\epsilon_{ijk}m_k .
\]
把空间相位在源尺度上展开到二阶，\(\ee^{-\ii k\vb n\cdot\vb x}=1-\ii k(\vb n\cdot\vb x)-\frac{k^2}{2}(\vb n\cdot\vb x)^2+\cdots\)，代入上述矩。零阶项 \(\int J_i\dd^3x=\dot p_i\) 经时间 Fourier 变换即电偶极 \(-\ii\omega\widetilde p_i\)。一阶项 \(-\ii k n_j\int x_j\widetilde J_i\dd^3x\) 中，\(-\epsilon_{ijk}m_k\) 部分给出磁偶极 \(\ii k\,\vb n\times\widetilde{\vb m}\)；\(\frac12\dot I_{ij}\) 部分经 \(\dot I_{ij}\to-\ii\omega\widetilde I_{ij}\) 后给出四极 \(-\frac{k\omega}{2}n_j\widetilde I_{ij}\)。

四极只有无迹部分存活。由 \(\mathcal Q^{\rm E}_{ij}=3I_{ij}-\delta_{ij}I_{kk}\) 得 \(I_{ij}=\frac13(\mathcal Q^{\rm E}_{ij}+\delta_{ij}I_{kk})\)，而横向投影 \(\Pi_{\perp,il}=\delta_{il}-n_in_l\) 满足 \(\Pi_\perp\cdot\vb n=0\)，故迹项 \(\delta_{ij}I_{kk}\) 在投影下整项消失：
\[
 \Pi_\perp\cdot(\vb I\cdot\vb n)
 =\frac13\Pi_\perp\cdot(\mathcal Q^{\rm E}\cdot\vb n).
\]
于是四极振幅 \(-\frac{k\omega}{2}\cdot\frac13=-\frac{k\omega}{6}\)。合并偶极、磁偶极与四极，并仍保留 \(O(\kappa_{\rm src}^2)\) 余项，便有
\[
 \vb\Pi_\perp\cdot\boldsymbol{\mathcal J}
 =-\ii\omega\vb\Pi_\perp\cdot\widetilde{\vb p}
 +\ii k\,\vb n\times\widetilde{\vb m}
 -\frac{k\omega}{6}\vb\Pi_\perp\cdot
 (\widetilde{\vb{\mathcal Q}}^{\rm E}\cdot\vb n)
 +O(\kappa_{\rm src}^2),
\]
即\eqnrefs{eq:phys-rr-multipole-current-amplitude}到所示阶次。

完整球面积分使用
\[
 \int n_in_j\dd\Omega=\frac{4\pi}{3}\delta_{ij},\qquad
 \int n_in_jn_kn_l\dd\Omega
 =\frac{4\pi}{15}
 (\delta_{ij}\delta_{kl}+\delta_{ik}\delta_{jl}+\delta_{il}\delta_{jk}).
\]
于是电偶极和磁偶极的横向角因子均为 \(8\pi/3\)。对对称无迹四极 \(\mathcal Q^{\rm E}_{ij}\)（以下记作 \(Q_{ij}\)），先把被积函数展开
\[
 |\Pi_\perp\cdot(\vb Q\cdot\vb n)|^2
 =Q_{ij}Q_{ik}n_jn_k-(\vb n\cdot\vb Q\cdot\vb n)^2,
\]
其中用了 \(n_in_i=1\) 与 \(Q_{ij}n_in_j=\vb n\cdot\vb Q\cdot\vb n\)。代入二阶、四阶角矩，
\[
 \int n_jn_k\dd\Omega=\frac{4\pi}{3}\delta_{jk},\qquad
 \int n_jn_kn_ln_m\dd\Omega=\frac{4\pi}{15}(\delta_{jk}\delta_{ln}+\delta_{jl}\delta_{kn}+\delta_{jn}\delta_{kl}).
\]
第一项为 \(\int Q_{ij}Q_{ik}n_jn_k=\frac{4\pi}{3}Q_{ij}Q_{ij}\)；第二项因 \(Q_{jj}=0\)，\(\int(\vb n\cdot\vb Q\cdot\vb n)^2=\frac{4\pi}{15}(Q_{jl}Q_{jl}+Q_{jj}Q_{kk}+Q_{jl}Q_{lj})=\frac{8\pi}{15}Q_{ij}Q_{ij}\)。相减得
\[
 \int\dd\Omega\,|\vb\Pi_\perp\cdot(\vb{\mathcal Q}^{\rm E}\cdot\vb n)|^2
 =\left(\frac{4\pi}{3}-\frac{8\pi}{15}\right)\mathcal Q^{\rm E}_{ij}(\mathcal Q^{\rm E}_{ij})^*
 =\frac{4\pi}{5}\mathcal Q^{\rm E}_{ij}(\mathcal Q^{\rm E}_{ij})^*.
\]
把这些结果代入精确母式，并使用
\(\mu_0\varepsilon_0c^2=1\)，得到
\eqnrefs{eq:phys-rr-e1-spectrum,eq:phys-rr-m1-spectrum,eq:phys-rr-e2-spectrum}。
更高多极来自相位展开的下一阶，余项仍由
\(\kappa_{\rm src}\) 而不是时间软参数控制。
\end{proof}
多极展开的物理后果集中在选择定则与宇称上。电 $2^L$ 极辐射 $E_L$ 带宇称 $(-1)^L$，磁多极 $M_L$ 带宇称 $(-1)^{L+1}$；故电偶极 $E1$（$L=1$）连接宇称相反的态，而 $M1$、$E2$ 连接宇称相同的态，二者不可混淆。源尺度 $a_{\rm src}$ 满足 $ka_{\rm src}\ll1$ 时，第 $L$ 阶多极功率相对 $E1$ 被压低 $(ka_{\rm src})^{2(L-1)}$：偶极主导、更高阶为禁戒。这解释了原子光谱中电偶极线远强于磁偶极与四极线——后两者只能经 $M1/E2$ 或更高阶跃迁，且自带 $(ka)^{2L}$ 的幂次抑制。

任意频率使用\eqnrefs{eq:phys-rr-exact-radiation-spectrum-current}；点世界线使用\eqnrefs{eq:phys-rr-exact-point-spectrum}；周期源使用\eqnrefs{eq:phys-rr-exact-periodic-line-power}。只有 \(\kappa_{\rm src}\ll1\) 时才能采用多极式\eqnrefs{eq:phys-rr-multipole-current-amplitude}，只有 \(\epsilon_{\rm soft}\ll1\) 时才能采用软因子\eqnrefs{eq:phys-rr-leading-soft-factor}。若轨道也要由动力学决定，就必须把外场与自场同时放入运动方程。
